\chapter{Implementación}
\label{cap:implementacion}

En este capítulo se detallan los aspectos más relevantes de la implementación del simulador en C++17, describiendo cómo las abstracciones matemáticas se han traducido a código, prestando atención a la gestión de memoria y las optimizaciones algorítmicas.

\section{El Cuerpo de Galois}
La implementación del cuerpo ($GF(2^m)$) es el componente más crítico del sistema, ya que todas las operaciones de codificación y decodificación dependen de su eficiencia.

\subsection{Tablas}
Se ha implementado la clase \verb|GaloisField| utilizando tablas de búsqueda para la multiplicación y división. Una decisión clave fue el uso de miembros estáticos para las tabla de logaritmos y exponenciales:

\begin{itemize}
    \item \verb|static std::vector<uint16_t> log_table|: Almacena el logaritmo en base $\alpha$ de cada elemento.
    \item \verb|static std::vector<uint16_t> exp_table|: Almacena las potencias de $\alpha$.
\end{itemize}
\begin{table}[h]
\centering
\begin{tabular}{|c|c|c|c|c|c|}
\hline
\textbf{Índice (i)} & 0 & 1 & 2 & \dots & $2^m-2$ \\ \hline
\textbf{exp\_table[i]} & $\alpha^0$ & $\alpha^1$ & $\alpha^2$ & \dots & $\alpha^{2^m-2}$ \\ \hline
\textbf{log\_table[val]} & $i=0$ & $i=1$ & $i=2$ & \dots & $i=2^m-2$ \\ \hline
\end{tabular}
\caption{Mapeo de las tablas de búsqueda en $GF(2^m)$ utilizando \texttt{uint16\_t}.}
\label{tab:tablas_gf}
\end{table}
El uso de \verb|static| asegura que, aunque el simulador cree múltiples instancias de la clase, los datos se comparten de forma global entre instancias, minimizando el consumo de memoria. La inicialización se realiza de forma perezosa, calculando las tablas únicamente cuando se detecta un cambio en el parámetro $m$ o en el polinomio primitivo. La elección de \verb|uint16_t| para las tablas es estratégica. Con $m=15$, cada tabla posee $2^{15}$ entradas, ocupando un total de 128 KB entre ambas, acelerando drásticamente las operaciones aritméticas repetitivas en el decodificador.

\section{Aritmética de Polinomios}
La clase \texttt{Polynomial} gestiona los coeficientes mediante un \texttt{std::vector<uint16\_t>}. Para mantener la integridad de las operaciones algebraicas sin comprometer el rendimiento, se han implementado las siguientes estrategias:

\subsection{Gestión de Grado y el Método \texttt{trim()}}
El método \texttt{trim()} es el encargado de ajustar el tamaño del vector eliminando coeficientes nulos de mayor grado. Para evitar una sobrecarga computacional innecesaria (especialmente en procesos iterativos), el diseño centraliza la ejecución de este método en los puntos de creación o modificación del objeto (constructor y \texttt{setCoef}). De esta forma, cualquier operación que genere un nuevo polinomio garantiza un estado consistente de forma atómica, eliminando la necesidad de comprobaciones redundantes durante el procesamiento de datos.

\subsection{Evaluación mediante el Algoritmo de Horner}
Para funciones críticas como el cálculo de síndromes, se emplea el método de Horner dentro de la función \texttt{evaluate()}. Esta técnica transforma la evaluación de un polinomio de grado $n$ en una secuencia de $n$ multiplicaciones y sumas, optimizando el uso de registros de CPU y evitando el cálculo explícito de potencias.

\section{Codificación Mediante \texttt{PackedBit}}
La codificación sistemática requiere dividir el mensaje extendido por el polinomio generador. Para acelerar este proceso, se ha implementado la estructura \texttt{PackedBit}.

\subsection{Procesamiento en Bloque}
En lugar de procesar bit a bit, \texttt{PackedBit} empaqueta la información en bloques de 64 bits (\texttt{uint64\_t}). El método \texttt{xorShifted} optimiza la división polinómica al evitar desplazamientos de bit individuales. El algoritmo calcula el desplazamiento de bloque (\textit{block\_shift}) y de bit (\textit{bit\_shift}) necesarios para alinear el polinomio generador con el dividendo, aplicando la operación XOR directamente sobre los \texttt{uint64\_t}. Esto permite procesar hasta 64 bits en un único ciclo de instrucción del procesador.

\begin{lstlisting}[caption={Implementación optimizada de XOR con desplazamiento en bloques de 64 bits.}]
// Alineacion y XOR en bloque para division polinomica
for (size_t i = 0; i < gx.data.size(); ++i) {
    size_t dst = i + block_shift;
    if (dst >= data.size()) break;
    
    uint64_t word = gx.data[i];
    
    // Gestion de acarreo (carry) para bits desplazados
    next_carry = (bit_shift == 0) ? 0 : (word >> (64 - bit_shift));
    
    // Operacioon XOR sobre el bloque de 64 bits completo
    data[dst] ^= (word << bit_shift) | carry;
    carry = next_carry;
}
\end{lstlisting}

\section{Implementación del Decodificador BCH}
La decodificación es la fase de mayor complejidad. Se ha optimizado mediante tres técnicas:

\subsection{Simetría de Síndromes}
Se aprovecha la propiedad de los campos binarios donde los síndromes de índice par satisfacen la relación $S_{2i} = (S_i)^2$. El simulador solo evalúa explícitamente los índices impares mediante Horner, calculando los pares mediante multiplicaciones directas en el campo. Esto reduce a la mitad la latencia de la etapa de detección de errores.

\subsection{Optimización de Berlekamp-Massey}
Para minimizar el \textit{overhead} de gestión de memoria dinámica, el algoritmo de Berlekamp-Massey se ha implementado operando directamente sobre vectores de coeficientes crudos. Se ha evitado el uso de la abstracción \texttt{Polynomial} en el bucle interno para prevenir la creación y destrucción constante de objetos temporales, lo cual reduciría la velocidad de simulación en escenarios con alta tasa de error.

\subsection{Búsqueda de Chien}
La localización de errores se realiza mediante una búsqueda de Chien. Se utilizan registros auxiliares (\texttt{reg}) y factores de actualización (\texttt{factors}) para evaluar el polinomio localizador en cada posición del bloque de forma incremental. Esta implementación garantiza una complejidad temporal lineal respecto al tamaño de bloque $n$.

\begin{figure}[h]
    \centering
    \begin{tikzpicture}[node distance=1.2cm, auto, >=latex', thick]
        \node (input) [draw, rectangle] {Vector Recibido $r(x)$};
        \node (syn) [draw, rectangle, below=of input] {Cálculo de Síndromes ($S_i$)};
        \node (bm) [draw, rectangle, below=of syn] {Algoritmo de Berlekamp-Massey};
        \node (elp) [draw, node distance=0.6cm, below=of bm, font=\footnotesize] {Polinomio Localizador $\sigma(x)$};
        \node (chien) [draw, rectangle, below=of elp] {Búsqueda de Chien};
        \node (corr) [draw, rectangle, below=of chien] {Corrección de Errores};
        
        \draw[->] (input) -- (syn);
        \draw[->] (syn) -- (bm);
        \draw[->] (bm) -- (elp);
        \draw[->] (elp) -- (chien);
        \draw[->] (chien) -- (corr);
    \end{tikzpicture}
    \caption{Etapas del proceso de decodificación BCH implementado.}
    \label{fig:decodificador_pasos}
\end{figure}

\section{Canal de Comunicaciones}
La clase \texttt{Channel} implementa la inyección de ruido para los modelos BSC y AWGN.

\subsection{Generación de Ruido Aleatorio}
Para las simulaciones de Monte Carlo, se utiliza el motor \texttt{std::mt19937} inicializado con \texttt{std::random\_device}. Esta elección garantiza un periodo de aleatoriedad lo suficientemente largo para evitar sesgos en el cálculo de la tasa de error de bit (BER) y la tasa de error de trama (FER). En el caso del canal AWGN, se emplea una modulación BPSK y una distribución normal, seguida de una etapa de decisión (\textit{hard decision}) para recuperar la secuencia binaria. Se ha preferido \texttt{std::mt19937} sobre \texttt{rand()} porque su enorme periodo ($2^{19937}-1$) ayuda a generar secuencias pseudoaleatorias de mayor calidad, evitando patrones en simulaciones largas. Para asegurar la validez estadística, el bucle de simulación monitoriza el número de errores detectados, deteniendo la ejecución al alcanzar un mínimo de 100 tramas erróneas por punto de Eb/N0, garantizando así una estimación fiable de la BER incluso en condiciones de bajo ruido.